% !TEX root = approx8.tex



%\section{Generalities on PCs and TPCs} \label{s:tpc-gnrl} 
%
%\subsection{Persistence modules} \label{sb:per-m} 
%
%  
%\INS \ say something about the ground filed/ring. Sometimes we'll use
%persistence modules $M^{\alpha}$ such that $M^{\alpha}$ is an
%$R$-module, e.g.~for $R=\Lambda_0$ etc. \medskip
%
%Let
%$M = \Bigl( \{M^{\alpha}\}_{\alpha \in \mathbb{R}}, \; i_{\alpha,
%  \beta}: M^{\alpha} \longrightarrow M^{\beta}, \; \forall \alpha \leq
%\beta \Bigr)$ be a persistence module. We will use the following
%shorthand notation throughout the paper:
%\begin{enumerate}
%\item For every $r \in \mathbb{R}_{\geq 0}$, $\alpha \in \mathbb{R}$
%  we write $s^{+}_r:M^{\alpha} \longrightarrow M^{\alpha+r}$ for the
%  map $i_{\alpha, \alpha+r}$.
%\item Let $u, v \in M^{\alpha}$ and $r \geq 0$. We write
%  $u \underset{r}{=}v$ to indicate that $s^{+}_r(u) = s^{+}_r(v)$.
%%\item For $\alpha \leq \beta$ we write
%%  $m_{\beta}: M^{\alpha} \longrightarrow M^{\beta}$ for the map
%%  $i_{\alpha, \beta}$.
%%\item For $u \in M^{\alpha}$, $v \in M^{\beta}$. We write
%%  $u \underset{\text{per}}{=} v$ to indicate that
%%  $m_{\gamma}(u) = m_{\gamma}(v)$, where
%%  $\gamma := \max \{\alpha, \beta \}$.
%\item For $u \in M^{\alpha}$ we denote 
%  $\ell(u) := \alpha$ and call $\ell(u)$ the persistence level of $u$.
%\end{enumerate}
%
%We will also denote by
%$$M_{\text{tot}} := \coprod_{\alpha \in \mathbb{R}} M^{\alpha}$$ the underlying set
%of $M$. Note that the persistence level $\ell$ can be viewed as a
%function $\ell: M_{\text{tot}} \longrightarrow \mathbb{R}$.  Finally
%we define $$M^{\infty} := \colim_{\alpha \to \infty} M^{\alpha},$$
%where the colimit (or direct limit) is taken with respect to the
%structural maps $i_{\alpha, \beta}$, $\alpha \leq \beta$, and call it
%the $\infty$-limit of $M$.
%
%\begin{rem}
%  Let $C$ be a filtered chain complex, with an increasing filtration
%  $C^{\alpha} \subset C^{\beta} \subset C$,
%  $\forall \; \alpha \leq \beta$. We will sometimes view $C$ as a
%  persistence module with
%  $i_{\alpha. \beta}: C^{\alpha} \longrightarrow C^{\beta}$ being the
%  inclusion maps. Now, let $u \in C^{\alpha}$ and consider the same
%  element $u$ but now viewed in $C^{\beta}$ for some $\beta > \alpha$.
%  For many practical purposes these two elements are viewed as one
%  element that belongs to $C$. However, from the persistence module
%  viewpoint $u \in C^{\alpha}$ and its image
%  $i_{\alpha, \beta}(u) \in C^{\beta}$ should be viewed as distinct
%  elements, both belonging to $C_{\text{tot}}$. In fact,
%  $C \neq C_{\text{tot}}$ and every element $u \in C$ appears
%  infinitely many times in $C_{\text{tot}}$. Finally, the function
%  $\ell: C_{\text{tot}} \longrightarrow \mathbb{R}$ should not be
%  confused with the following filtration level measurement:
%  $$\ell_f : C \longrightarrow \mathbb{R} \cup \{\pm \infty\},
%  \quad \ell_{f} (c) : = \inf \{ \alpha \mid c \in C^{\alpha} \}$$
%  which is often used in the context of filtered chain complexes.
%\end{rem}
%
%
%\subsection{PC's} \label{sb:pc} {\em Persistence Categories} (PC's in
%short) are categories in which the morphisms between objects form
%persistence modules and the composition of morphisms is compatible
%with the persistence structures. The general theory of persistence
%categories is described in detail in~\cite[Section~2]{BCZ:tpc} and
%here we will only recall several key concepts and fix some notation.
%
%Unless otherwise stated, from now on we will always assume all the PC
%categories to be endowed with a shift functor (or more precisely a
%system of shift functors) $\Sigma = \{\Sigma^r\}_{r \in \mathbb{R}}$
%(see~\cite[Section~2.2.3]{BCZ:tpc} for the definition), and all the PC
%functors to commute with $\Sigma^r$ for every $r \in \mathbb{R}$.
%
%Let $\mathscr{C}$ be a PC category. Denote by $\mathscr{C}^0$ the
%$0$-level category associated with $\mathscr{C}$ and by
%$\mathscr{C}^{\infty}$ the limit (or $\infty$-level) category of
%$\mathscr{C}$. Isomorphisms in $\msc^0$ will be called also
%$0$-isomorphisms and sometimes denoted by $\cong_0$. Isomorphisms in
%$\msc^{\infty}$ will sometimes be denoted by $\cong_{\infty}$.
%
%Most of the time we will assume our PC's $\mathscr{C}$ to be graded in
%the sense that the persistence modules $\hom_{\mathscr{C}}(A,B)$ are
%$\mathbb{Z}$-graded. We will mostly use cohomological grading
%conventions and denote by $\hom_{\mathscr{C}}(A,B)^k$,
%$k \in \mathbb{Z}$ the degree-$k$ component of
%$\hom_{\mathscr{C}}(A,B)$. For $\alpha \in \mathbb{R}$ we denote by
%$\hom^{\alpha}_{\mathscr{C}}(A,B)^k$ the $\alpha$-persistence level of
%$\hom_{\mathscr{C}}(A,B)^k$. We will assume composition of morphisms
%to be degree-preserving and identity morphisms to lie in degree $0$.
%
%Our PC's $\mathscr{C}$ will be often endowed with a {\em translation
%  functor} $T: \mathscr{C} \longrightarrow \mathscr{C}$. This is a PC
%isomorphism functor with the property that for all
%$A, B \in \Ob(\mathscr{C})$, $k \in \mathbb{Z}$, we have PC-natural
%isomorphisms
%$$\hom_{\mathscr{C}}(TA,B)^k \cong \hom_{\mathscr{C}}(A,B)^{k-1},
%\quad \hom_{\mathscr{C}}(A,TB)^k \cong
%\hom_{\mathscr{C}}(A,B)^{k+1}.$$ As mentioned earlier, since $T$ is a
%PC functor, we will implicitly assume that $T$ commutes with the shift
%functors, i.e.~$T \circ \Sigma^{r} = \Sigma^r \circ T$,
%$\forall \, r \in \mathbb{R}$.
%
%
%
%
%\subsubsection{Acyclic objects} \label{sbsb:acyclic} Let $\mathscr{C}$
%be a PC, $A \in \Ob(\mathscr{C})$ and $r \geq 0$. We say that $A$ is
%$r$-acyclic if the morphism $\eta_r^A: \Sigma^r A \longrightarrow A$
%equals $0$. Here,
%$\eta^A_{r} \in \hom_{\mathscr{C}^0}(\Sigma^{r}A, A)$ is the standard
%map associated with the persistence structure of $\mathscr{C}$. An
%object $A$ is called acyclic if it is $r$-acyclic for some $r \geq
%0$. Equivalently, $A$ is acyclic if and only if $A$ is isomorphic to
%$0$ in the limit category $\mathscr{C}^{\infty}$.  We denote by
%$A\mathscr{C}$ the full subcategory of $\mathscr{C}$ whose objects are
%the acylic ones.
%
%
%\subsubsection{Pseudo metrics associated to
%  PC's} \label{sbsb:pc-metrics}
%
%Let $\mathscr{C}$ be a PC with a shift functor $\Sigma$. The
%collection of objects $\Ob(\mathscr{C})$ carries a pseudo-metric
%$\dint$, called the interleaving distance, which is defined as
%follows.  For $X, Y \in \Ob(\mathscr{C})$ define:
%\begin{equation} \label{eq:d-int}
%  \begin{aligned}
%    \dint(X, Y) = \inf \Bigl\{r \geq 0 & \bigm| \exists \, \varphi:
%    \Sigma^rX \longrightarrow Y, \, \exists \, \psi: \Sigma^rY
%    \longrightarrow X \\
%    & \;\;\; \text{such that} \; \psi \circ \Sigma^r \varphi =
%    \eta_{2r}^X, \; \varphi \circ \Sigma^r \psi = \eta_{2r}^Y \Bigr\}.
%  \end{aligned}
%\end{equation}
%Here, $\eta^X_{2r} \in \hom_{\mathscr{C}^0}(\Sigma^{2r}X, X)$ are the
%standard maps associated with the persistence structure of
%$\mathscr{C}$ and its shift functor and similarly for $\eta^Y_{2r}$.
%
%Note that $\dint$ is in general not a genuine metric but only a
%pseudo-metric and moreover it may have infinite values. Indeed, if $X$
%and $Y$ are isomorphic in the subcategory $\mathscr{C}^0$, then
%$\dint(X, Y) = 0$. Similarly, if $X$ and $Y$ are not isomorphic in
%$\mathscr{C}^{\infty}$ then $\dint(X,Y) = \infty$. This is compatible
%with the convention that $\inf \emptyset = \infty$ which we use
%in~\eqref{eq:d-int}.
%
%There is another variant of $\dint$ which measures how far an object
%$R$ is from being a retract of another object $X$. More precisely, for
%$R, X \in \Ob(\mathscr{C})$ we define
%\begin{equation} \label{eq:d-rint} \dret(R,X) = \inf \Bigl\{r \geq 0
%  \bigm| \exists \, \varphi: \Sigma^rR \longrightarrow X, \, \exists
%  \, \psi: \Sigma^rX
%  \longrightarrow R \\
%  \;\; \text{such that} \; \psi \circ \Sigma^r \varphi = \eta_{2r}^R
%  \Bigr\}.
%\end{equation}
%In contrast to $\dint$, the measurement $\dret(-,-)$ is in general not
%symmetric.
%
%In what follows we will use the following notation. Let $Z$ be a set
%endowed with a (not necessarily symmetric) pseudo-metric $d$ and
%$X, Y \subset Z$. We denote
%$$\vec{d}(X,Y) := \sup_{x \in X} \inf_{y \in Y} d(x,y)$$ the longest distance
%between $X$ and $Y$. (The arrow above the $d$ indicates that this is
%not symmetric in $(X,Y)$.) We will sometimes use this measurement in
%the case of $d=\dint$ and denote it $\vec{d}_{\text{int}}$. Similarly
%we can also define also $\vec{d}_{\text{r-int}}$ (although $\dret$ is
%not a symmetric pseudo-metric to start with).
%
%To end this section, we list a few basic metric properties of PC
%functors. In the following lemma the same notation $\dint$ is used for
%the inteleaving distance in various PC categories.
%\begin{lem} \label{l:d-func} Let $\msc', \msc''$ be PC's and
%  $\msf: \msc' \longrightarrow \msc''$ a PC-functor.
%  \begin{enumerate}
%  \item For every $A, B \in \Ob(\msc')$ we have
%    $\dint(\msf A, \msf B) \leq \dint(A,B)$.
%  \item If there exists a PC functor
%    $\msg: \msc'' \longrightarrow \msc'$ such that
%    $\dint(\msg \circ \msf, \id_{\msc'}) \leq s$ then for every
%    $A, B \in \Ob(\msc')$ we have
%    $|\dint(\msf A, \msf B) - \dint (A,B)| \leq 2s$.
%  \item If there is a PC-functor $\msh: \msc'' \longrightarrow \msc'$
%    such that $\dint(\msf \circ \msh, \id_{\msc''}) \leq r$ then
%    $$\vec{d}_{\text{int}}(\msc'', \textnormal{image\,}\msf) \leq r.$$
%  \end{enumerate}
%\end{lem}
%
%\INS: explain the lemma for the case of $\dret$.

\subsection{Conventions for PC's and filtered $A_{\infty}$
  categories}\label{subsec:A_{infty}}

\subsubsection{Filtered $A_{\infty}$-categories} \label{sbsb:fai} Most
of the PC's in this paper arise as persistence homological categories
of filtered $A_{\infty}$-categories. Unless otherwise stated we will
always assume such categories, as well as $A_{\infty}$-functors
between them, to be strictly unital and endowed with a shift and
translation functors. See~\S\ref{a:fil-ai} for the precise
definitions.

Let $\mathcal{A}$ be a filtered $A_{\infty}$-category.  We will
sometimes abbreviate $\mathcal{A}(X,Y) :=
\hom_{\mathcal{A}}(X,Y)$. For $X, Y \in \Ob(\mathcal{A})$,
$\alpha \in \mathbb{R}$, we denote by
$\hom_{\mathcal{A}}^{\alpha}(X,Y) \subset \hom_{\mathcal{A}}(X,Y)$, or
sometimes by $\mathcal{A}^{\alpha}(X,Y)$ for short, the
$\alpha$-filtration level of $\hom_{\mathcal{A}}(X,Y)$. Whenever we
need to combine these with (cohomological) grading we will write
$\hom_{\mathcal{A}}^{\alpha}(X,Y)^i$ (or
$\mathcal{A}^{\alpha}(X,Y)^i$) for the degree-$i$ component of
$\hom_{\mathcal{A}}^{\alpha}(X,Y)$. The notation in case of
homological grading is similar, by making the degree $i$ a subscript.

Next, we denote by $\mathcal{A}^0$ the $0$-level category of
$\mathcal{A}$, namely the (unfiltered) $A_{\infty}$-category with the
same objects as $\mathcal{A}$ and
$\hom_{\mathcal{A}^0}(X,Y) = \hom_{\mathcal{A}}^0(X,Y)$.

Passing to homology, we write $PH(\mathcal{A})$ for the persistence
homological category of $\mathcal{A}$. The shift and translation
functors induce respective functors on $PH(\mathcal{A})$.

As in the case of PC's, we say that an object $A \in \mathcal{A}$ is
$r$-acyclic if $A$ is $r$-acyclic in the persistence category
$PH(\mathcal{A})$, and similarly for acyclic objects (without
reference to a specific $r$). The full $A_{\infty}$-subcategory of
acyclic objects will be denoted by $A\mathcal{A}$.

Finally, we import the interleaving (and r-interleaving)
  pseudo-metrics also to the realm of $A_{\infty}$-categories
$\mathcal{A}$, by setting both of them to be equal to the respective
distances (as defined in~\S\ref{sbsb:pc-metrics}) in the persistence
homological category $PH(\mathcal{A})$.


\subsubsection{Various completions} \label{sbsb:comp} Let $\msc$ be a
PC. Given subsets $S_1, S_2, S_3 \subset \Ob(\mathscr{C})$ we define
their completions with respect to shifts, translations and
$0$-isomorphisms, respectively, as follows:
\begin{equation}
  \begin{aligned}
    & S_1^{\Sigma} := \{ \Sigma^r(A) \mid A \in S_1, \, r \in
    \mathbb{R}\}, \quad
    S_2^{T} := \{T^j(A) \mid A \in S_2, \, j \in \mathbb{Z}\}, \\
    & S_3^{\ziso} := \{ A \in \Ob(\mathscr{C}) \mid A \text{ is
      isomorphic in } \mathscr{C}^0 \text{ to an object from } S_3\}.
  \end{aligned}
\end{equation}
We will also need the completion of a subset
$S \subset \Ob(\mathscr{C})$ with respect to several of the above
procedures, for example $S^{\Sigma, T} := (S^{\Sigma})^T$ etc.

Another important completion is with respect to the interleaving
distance. Given a subset $S_4 \subset \Ob(\msc)$ we define
$$S_4^{\text{int}} :=
\{ A \in \Ob(\mathscr{C}) \mid \dint(A, B) < \infty \text{ for some }
B \in S_4\}.$$ To simplify the notation, for $S \subset \Ob{\msc}$ we
will write $S^{c} := (S^{T})^{\text{int}}$.
%$$S^{c,0} := S^{\Sigma, T, \ziso}, \quad S^{c} := (S^{T})^{\text{int}}.$$
Note that $S^c$ is automatically complete with respect to both shifts
and translations.
% completing with respect to the interleaving distance already
% includes the completion with respect to shifts and to
% $\msc^0$-isomorphisms, hence we can omit them when defining $S^c$.
One can also define $S^{c,\text{r-int}}$ in a similar way to $S^c$,
but with $\dint$ replaced by $\dret$.

If $\mathscr{D} \subset \mathscr{C}$ is a full sub-PC we define its
various completions, e.g.~$\mathscr{D}^{\Sigma}$, $\mathscr{D}^{T}$,
$\mathscr{D}^{(\text{iso}, 0)}$, $\mathscr{D}^{c}$ etc.~by taking the
full sub-PC's of $\mathscr{C}$ with the respective completed sets of
objects $\Ob(\mathscr{D})^{\Sigma}$, $\Ob(\mathscr{D})^{T}$,
$\Ob(\mathscr{D})^{\ziso}$, $\Ob(\mathscr{D})^{c}$ etc. In the case of
$A_{\infty}$-categories $\mathcal{A}$, we define analogous
completions, by completing the subsets of objects using the respective
structures in $PH(\mathcal{A})$.

%\subsection{TPC's} \label{sb:tpc} A {\em Triangulated Persistence
%  Category} (TPC in short) $\mathscr{C}$ is a PC with an additional
%structure which makes its $0$-level category $\mathscr{C}^0$ a
%triangulated category. The precise definition of TPC involves several
%other axioms and we refer the reader to~\cite{BCZ:tpc} for a detailed
%foundation of the subject.
%
%Let $\msc$ be a TPC and $S \subset \Ob(\msc)$. We write
%$\langle S \rangle^{\Delta}$ for the minimal full sub-TPC of $\msc$
%which contains all the objects of $S$ and is closed under
%$0$-isomorphisms. (Note that in particular the shift and translation
%functors of $\msc$ restrict to respective functors on
%$\langle S \rangle^{\Delta}$.) We denote by
%$(\langle S \rangle^{\Delta})^0$ the $0$-level category of the TPC
%$\langle S \rangle^{\Delta}$.

\subsection{Persistence derived categories associated with a filtered
  $A_{\infty}$-category} \label{sb:pdc} We will present below three
variants of TPC's that can be viewed each as a generalization of the
derived category of $\mathcal{A}$ to the realm of persistence
categories. However before we go into this we need a quick preparation
about $A_{\infty}$-functors in the filtered context.

\subsubsection{Filtered functors and functors with linear
  deviation} \label{sbsb:fil-func}

Let $\mathcal{A}$, $\mathcal{B}$ be two filtered
  $A_{\infty}$-categories. A filtered $A_{\infty}$-functor
  $\mathcal{F}: \mathcal{A} \longrightarrow \mathcal{B}$ is an
  $A_{\infty}$-functor whose action on morphisms (in all orders)
  preserves the filtration levels. More specifically, for every
  $d \geq 1$, $X_0, \ldots, X_d \in \Ob(\mathcal{A})$ and
  $\alpha \in \mathbb{R}$ the $d$'th order term
  $\mathcal{F}_d: \mathcal{A}(X_0, \ldots, X_d) \longrightarrow
  \mathcal{B}(\mathcal{F}X_0, \mathcal{F}X_d)$ of $\mathcal{F}$
  satisfies:
  $$\mathcal{F}_d \bigl(\mathcal{A}^{\alpha}(X_0, \ldots, X_d)\bigr) \subset
  \mathcal{B}^{\alpha}(\mathcal{F}X_0, \mathcal{F}X_d).$$ Unless
  otherwise stated, in what follows we will implicitly assume our
  $A_{\infty}$-functors to be strictly unital.

 Motivated by examples coming from Fukaya categories we will
  also need the concept of $A_{\infty}$-functors with linear
  deviation. These are defined as follows. Given a tuple
  $\vec{X} = (X_0, \ldots, X_d)$ of objects from $\mathcal{A}$ we
  define its reduced tuple
  $\vec{X}_R := (X_{i_0}, \ldots, X_{i_{d_R}})$ by ommitting from
  $\vec{X}$ subsequent (in the cyclic order) objects that are equal up
  to a shift. The objects forming $\vec{X}_R$ are well defined only up
  to shifts, but the length of $\vec{X}_R$,  $0 \leq d_R \leq d$, is well defined. We call it the reduced length
  of $\vec{X}$ and denote it by $d_R$ or $d_R(\vec{X})$ whenever we
  want to emphasize its dependence on $\vec{X}$. Note that in case
  every two consecutive objects in $\vec{X}$ are different (up to
  shifts) then $d_R(\vec{X}) = d$. At the other extreme, if all the
  objects in $\vec{X}$ are equal up to shifts, then
  $d_R(\vec{X}) = 0$.  

An $A_{\infty}$-functor
  $\mathcal{F}: \mathcal{A} \longrightarrow \mathcal{B}$ is said to
  have linear deviation rate $s \geq 0$ if for every $d \geq 1$, every
  tuple of objects $\vec{X} = (X_0, \ldots, X_d)$ from
  $\Ob(\mathcal{A})$ and every $\alpha \in \mathbb{R}$ we have:
$$\mathcal{F}_d \bigl(\mathcal{A}^{\alpha}(X_0, \ldots, X_d)\bigr)
\subset \mathcal{B}^{\alpha + d_R(\vec{X}) s}(\mathcal{F}X_0,
\mathcal{F}X_d).$$ We will often refer to such functors as LD (Linear
Deviation) functors.


\subsubsection{Filtered modules} \label{sbsb:fmod} Let
$\mathcal{A}$ be a filtered $A_{\infty}$-category. Denote by $F
\md_{\mathcal{A}}$ the category of filtered strictly unital (left)
$\mathcal{A}$-modules. This is a filtered 
$A_{\infty}$-category. The Yoneda embedding
$$\mathcal{Y}: \mathcal{A} \longrightarrow F\md_{\mathcal{A}}$$
is a filtered strictly unital $A_{\infty}$-functor which is
homologically full and faithful.  Denote by
$\mathcal{Y}(\mathcal{A}) \subset F \md_{\mathcal{A}}$ the image
category of $\mathcal{A}$ under $\mathcal{Y}$. Denote by
$\mathcal{Y}(\mathcal{A})^{\Delta} \subset F \md_{\mathcal{A}}$ the
triangulated completion of $\mathcal{Y}(\mathcal{A})^{\Delta}$,
namely the minimal full subcategory of $F \md_{\mathcal{A}}$ which is
closed under mapping cones over morphisms with filtration level
$\leq 0$ and quasi-isomorphisms of filtration level $\leq 0$
(these are maps of non-positive filtration level that induce an isomorphism in the $0$-level homological category). Note
that due to our assumptions on the filtered $A_{\infty}$ category $\mathcal{A}$, the category
$\mathcal{Y}(\mathcal{A})$ is closed under shifts and translations,
and therefore the same holds for $\mathcal{Y}(\mathcal{A})^{\Delta}$.
The construction of $\mathcal{Y}(\mathcal{A})^{\Delta}$ can be
explicitly carried out by iteratively taking mapping cones over
morphisms of filtration level $\leq 0$ and adding at each stage
$0$-quasi-isomorphic modules. The resulting
$A_{\infty}$-category $\mathcal{Y}(\mathcal{A})^{\Delta}$ is filtered
and strictly unital $A_{\infty}$ and carries a translation and a shift
functor. The persistence-homological category of
$\mathcal{Y}(\mathcal{A})^{\Delta}$
$$PD(\mathcal{A}) := PH(\mathcal{Y}(\mathcal{A})^{\Delta})$$
is a TPC which we call the {\em persistence derived} category of
$\mathcal{A}$.

Let $\mathcal{A}$, $\mathcal{B}$ be two filtered
$A_{\infty}$-categories. Let
  $\mathcal{F}: \mathcal{A} \longrightarrow \mathcal{B}$ be an
  $A_{\infty}$-functor with a given linear deviation rate
  (see~\S\ref{sbsb:fil-func} and~\S\ref{a:func-LD}). Then the
  push-forward operation from~\S\ref{ap:pshf} gives rise to a {\em
    filtered} $A_{\infty}$-functor
  $\mathcal{F}_*: F\md_{\mathcal{A}} \longrightarrow
  F\md_{\mathcal{B}}$.  More importantly for us, it sends
$\mathcal{Y}(\mathcal{A})^{\Delta}$ to
$\mathcal{Y}(\mathcal{\mathcal{B}})^{\Delta}$ hence it gives rise to a
filtered $A_{\infty}$-functor
$$\mathcal{F}_* : \mathcal{Y}(\mathcal{\mathcal{A}})^{\Delta}
\longrightarrow \mathcal{Y}(\mathcal{\mathcal{B}})^{\Delta}.$$ Passing
to homology we obtain a TPC functor
$$PD(\mathcal{F}): PD(\mathcal{A}) \longrightarrow PD(\mathcal{B}).$$

\subsubsection{A larger derived category} \label{sbsb:fmod-2} As in
the previous sections, let $\mathcal{A}$ be a filtered
$A_{\infty}$-category. We will now construct a somewhat larger version
of the category $PD(\mathcal{A})$ from~\S\ref{sbsb:fmod} by adding
also acyclic modules.

Denote by $AF\md_{\mathcal{A}} \subset F\md_{\mathcal{A}}$ the full
$A_{\infty}$-subcategory of acyclic modules (see~\S\ref{sbsb:fai} for
the definition of the acyclic subcategory). Consider now the minimal
$A_{\infty}$ full subcategory of $F \md_{\mathcal{A}}$ which contains
both the image of the filtered Yoneda embedding
$\mathcal{Y}(\mathcal{A})$ as well as $AF\md_{\mathcal{A}}$ and is
closed under mapping cones over morphisms with filtration level
$\leq 0$ and quasi-isomorphisms of filtration level $\leq 0$.  Denote
this filtered $A_{\infty}$-category by
$\langle \mathcal{Y}(\mathcal{A}), AF\md_{\mathcal{A}}
\rangle^{\Delta}$. Define now the complete persistence derived
category of $\mathcal{A}$ to be the persistence homology category of
the latter, namely:
$$PD^{c}(\mathcal{A}) := PH \Bigl( \langle \mathcal{Y}(\mathcal{A}),
AF\md_{\mathcal{A}}\rangle^{\Delta} \Bigr).$$ The discussion
from~\S\ref{sbsb:fmod} on functors carries over to this case. More
specifically, if $\mathcal{A}$, $\mathcal{B}$ are filtered
$A_{\infty}$-categories and
$\mathcal{F}: \mathcal{A} \longrightarrow \mathcal{B}$ is an
$A_{\infty}$-functor with linear deviation then we obtain a TPC
functor:
$$PD^c(\mathcal{F}): PD^c(\mathcal{A})
\longrightarrow PD^c(\mathcal{B})$$ induced by $\mathcal{F}$.


There is an alternative description of $PD^c(\mathcal{A})$, by
completing $PD(\mathcal{A})$ with respect to the interleaving distance, or even by applying  the same
type of completion to $\mathcal{Y}(\mathcal{A})^{\Delta}$ and then
passing to persistence homology. The result turns out to be the same:
\begin{lem} \label{l:PDc}
  $PD^{c}(\mathcal{A}) = \bigl(PD(\mathcal{A})\bigr)^{c} = PH \Bigl(
  (\mathcal{Y}(\mathcal{A})^{\Delta})^{c} \Bigr)$.
\end{lem}

\subsubsection{$r$-isomorphisms} \label{sbsb:r-iso} In the context of
TPC's there is an important notion of $r$-isomorphism that leads to
interesting measurements and can also be used to arrive at the same
completion $PD^c(\mathcal{A})$ described in~\S\ref{sbsb:fmod-2}.

Let $\mathscr{C}$ be a TPC, $A,B \in \Ob(\mathscr{C})$,
$\phi: A \longrightarrow B$ a morphism in $\hom_{\mathscr{C}^0}(A,B)$
and $r \geq 0$. We say that $\phi$ is an $r$-isomorphism if $\phi$ can
be completed to an exact triangle in the (triangulated) category
$\mathscr{C}^0$
$$A \xrightarrow{\ \phi \ } B \xrightarrow{\ \ \ }
K \xrightarrow{\ \ \ } TA$$ with $K \in \Ob(\mathscr{C})$ $r$-acyclic
(see~\S\ref{sbsb:acyclic}). For example, the standard morphism
$\eta^A_r : \Sigma^r A \longrightarrow A$ is always an $r$-isomorphism
(this is in fact one of the axioms defining TPC's), but of course
there are many other examples. We refer the reader to~\cite{BCZ:tpc}
for more details on the concept of $r$-isomorphisms.

Completing with respect to $r$-isomorphisms is somewhat subtle because
of the following. The existence of an $r$-isomorphism
$\phi: A \longrightarrow B$ does not imply that there exists an
$r$-isomorphism $B \longrightarrow A$, hence this does not lead to a
symmetric relation on pairs of objects $(A,B)$. Nor is this
``relation'' transitive (if $\phi: A \longrightarrow B$ is an
$r_1$-isomorphism and $r_2$-isomorphism $\psi: B \longrightarrow C$ is
an $r_2$-isomorphism then $\psi \circ \phi: A \longrightarrow C$ is in
general only an $(r_1+r_2)$-isomorphism). However by the results
of~\cite{BCZ:tpc}, if $\phi: A \longrightarrow B$ is an
$r$-isomorphism then there exists a $2r$-isomorphism
$\Sigma^r B \longrightarrow A$. This leads to the following
definition: two objects $A, B \in \Ob(\msc)$ are called {\em almost
  isomorphic} if there is an $r\geq 0$ and $s \in \mathbb{R}$ and an
$r$-isomorphism $\phi: A \longrightarrow \Sigma^s B$. By the results
of~\cite{BCZ:tpc}, being ``almost isomorphic'' is an equivalence
relation. Moreover, we have:
\begin{lem} \label{l:a-iso} Let $\msc$ be a TPC and Let
  $A, B \in \Ob(\msc)$. Then $A$ is almost isomorphic to $B$ if and
  only if $\dint(A,B) < \infty$.
\end{lem}
It follows that completing $PD(\mathcal{A})$ with respect to almost
isomorphisms gives precisely the same category as $PD^c(\mathcal{A})$.

To end this discussion, let us also mention the following useful fact:
completing a TPC with respect to the interleaving distance always
yields a TPC. More precisely
\begin{lem} \label{l:cmp-TPC} Let $\mathscr{C}$ be a TPC and
  $\mathscr{D} \subset \mathscr{C}$ a full sub-TPC of
  $\mathscr{C}$. Then $\mathscr{D}^{c} \subset \mathscr{C}$ is also a
  sub-TPC.
\end{lem}

\subsubsection{Filtered twisted complexes} \label{sbsb:ftw} Let
$\mathcal{A}$ be a filtered $A_{\infty}$-category. There is another
variant of a TPC associated with $\mathcal{A}$ which is based on
twisted complexes (see~\cite[Chapter~I,
Section~3(l)]{Se:book-fukaya-categ} for unfiltered
$A_{\infty}$-categories and~\cite[Section~2.5]{BCZ:tpc} for the case
of filtered dg-categories. The case of filtered
$A_{\infty}$-categories is treated in detail in~\S\ref{ap:ftc}.

Roughly speaking, one defines first the additive enlargement
$\mathcal{A}^{\oplus}$ of $\mathcal{A}$ as
in~\cite{Se:book-fukaya-categ} (where it is denoted by
$\Sigma \mathcal{A}$, but in our case $\Sigma$ is already used for
other purposes), and defines a filtration structure on it by an
obvious extension of the filtrations on $\mathcal{A}$. Recall that, by
assumption, $\mathcal{A}$ is closed under shifts and translations.
Next one forms a new filtered $A_{\infty}$-category $F Tw \mathcal{A}$
whose objects are all twisted complexes $(X, q_X)$ over $\mathcal{A}$
with $X \in \Ob(\mathcal{A}^{\oplus})$ and differential
$q_X \in \hom^{0}_{\mathcal{A}^{\oplus}}(X,X)$ lying at filtration
level $\leq 0$. The morphisms in this category come from
$\mathcal{A}^{\oplus}$ and are thus filtered. One then defines the
$A_{\infty}$-operations on $FTw \mathcal{A}$ in the standard way and
the fact that $\mathcal{A}$ itself is filtered implies that
$FTw \mathcal{A}$ is filtered too. Finally note that since
$\mathcal{A}$ is strictly unital the same holds for
$F Tw \mathcal{A}$.

The filtered $A_{\infty}$-category $\mathcal{A}$ embeds into
$Tw \mathcal{A}$ in an obvious way, the embedding being a filtered
$A_{\infty}$-functor which is full and faithful (on the chain
level). Moreover, $Tw \mathcal{A}$ is pre-triangulated in the filtered
sense (which in particular means that it is closed under formation of
filtered mapping cones). We denote by $PH(F Tw \mathcal{A})$ the
persistence homological category of $F Tw \mathcal{A}$. Standard
arguments show that $PH(F Tw \mathcal{A})$ is a TPC with translation
and shift functors induced from those of $\mathcal{A}$.

An important property of $PH(F Tw \mathcal{A})$ is the following.  Let
$\mathcal{A}$, $\mathcal{B}$ be two filtered
$A_{\infty}$-categories. Let
$\mathcal{F}: \mathcal{A} \longrightarrow \mathcal{B}$ be a filtered
$A_{\infty}$-functor. Then there is a canonical extension of
$\mathcal{F}$ to a filtered, strictly unital, $A_{\infty}$-functor
$Tw \mathcal{F}: FTw \mathcal{A} \longrightarrow FTw \mathcal{B}$.
Moreover, $Tw \mathcal{F}$ induces a homological functor
$$PH(Tw \mathcal{F}): PH(F Tw \mathcal{A}) \longrightarrow PH(FTw
\mathcal{B})$$ which is a TPC functor. 